% 第5章 工作计划

\chapter{五、工作计划}

本课题按五个阶段组织实施，各阶段的工作内容、工作量、起讫日期与阶段成果形式如表 5.1 所示。

\begin{table}[htbp]
  \centering
  \caption{研究工作进度安排}
  \label{tab:plan}
  \zihao{-5}
  \begin{tabular}{p{0.6cm}p{4.7cm}p{1.8cm}p{2.1cm}p{3.2cm}}
    \hline
    序号 & 阶段及内容 & 工作量估计（时数） & 起讫日期 & 阶段成果形式 \\
    \hline
    1 & \textbf{第一阶段：理论准备与方案细化} & 约 200 小时 & 2025.10--2026.03 & 1. 完成文献综述报告 \\
      & 深度文献调研，完善``线定量 + 场成像''互校验研究框架 & & & 2. 提交详细的实验与仿真技术方案书 \\
      & 完成 QCL-TDLAS 与 burst 染料激光 NO-PLIF 实验方案及仿真模型详细设计 & & & \\
    \hline
    2 & \textbf{第二阶段：实验系统搭建与仿真预研} & 约 600 小时 & 2026.04--2026.08 & 1. 可稳定运行的实验平台 \\
      & 搭建中红外 QCL-TDLAS 平台：折返式长光程吸收池装调与有效光程标定、波长扫描与``吸收谱—温度—浓度''联合反演程序 & & & 2. 仿真预研报告，包含工况建议与光谱数据库筛选结果 \\
      & 搭建 burst 染料激光 NO-PLIF 平台 & & & \\
      & 光谱仿真确定实验工况与谱线 & & & \\
    \hline
    3 & \textbf{第三阶段：实验测量与数据采集} & 约 800 小时 & 2026.09--2027.03 & 1. 完整原始光谱与图像数据集 \\
      & 标准火焰上开展多工况 QCL-TDLAS 与 NO-PLIF 同步测量 & & & 2. 初步浓度/流场结构观测报告 \\
      & 完成重复性实验与``点—面''同步采集并在发动机尾气工况下做现场验证 & & & \\
    \hline
    4 & \textbf{第四阶段：数据处理、反演与模型构建} & 约 600 小时 & 2027.04--2027.08 & 1. NO 浓度/温度场分布图谱 \\
      & 完成``吸收谱—温度—浓度''联合反演 & & & 2. 互校验分析报告与测量不确定度模型 \\
      & 建立 TDLAS 路径积分浓度与 PLIF 二维浓度场的空间映射与权重构造 & & & \\
      & 实验—仿真对比（Cantera），偏差评估 & & & \\
    \hline
    5 & \textbf{第五阶段：论文撰写与成果凝练} & 约 300 小时 & 2027.09--2027.12 & 1. 提交完整的硕士学位论文 \\
      & 完成学位论文撰写、修改与润色 & & & 2. 发表学术论文 1--2 篇（或投稿） \\
      & 整理研究成果，准备答辩 & & & 3. 答辩 PPT 及相关材料 \\
    \hline
    \multicolumn{2}{l}{合计} & 约 2500 小时 & & \\
    \hline
  \end{tabular}
\end{table}
